\section{Expansão da Implementação do Modelo de \textit{Backlash}}

Equações do modelo de expansão:

\begin{equation}
	\begin{split}
		\delta(t) &= \left[ (x_1 - x_2) \sin\psi + (y_1 - y_2) \cos\psi + R_1\theta_{z1} + R_2\theta_{z2} \right] \cos\beta_h \\
		&\quad + \left[ (z_2 - z_1) + (R_1\theta_{x1} + R_2\theta_{x2}) \sin\psi + (R_1\theta_{y1} + R_2\theta_{y2}) \cos\psi \right] \sin\beta_h - e(t)
	\end{split}
	\label{eq:delta_3d}
\end{equation}

\begin{equation}
	b_t = b_0 + (R_1 + R_2) (\mathrm{inv}(\alpha) - \mathrm{inv}(\alpha_0)) \cos\beta_h
	\label{eq:bt_helicoidal}
\end{equation}

\begin{equation}
	\begin{aligned}
		\partial_{x_i}b_t &= (R_1 + R_2) \tan^2\alpha \cdot \partial_{x_i}\alpha \cdot \cos\beta_h \\
		\partial_{y_i}b_t &= (R_1 + R_2) \tan^2\alpha \cdot \partial_{y_i}\alpha \cdot \cos\beta_h
	\end{aligned}
	\label{eq:dbt_helicoidal}
\end{equation}

\text{Para as derivadas parciais do DTE (com } $\psi = \alpha - \beta $ \text{), definem-se os termos geométricos auxiliares:}

\begin{equation}
	\begin{aligned}
		\Delta_{\text{trans}} &= (x_1 - x_2) \cos\psi - (y_1 - y_2) \sin\psi \\
		\Delta_{\text{rot}} &= (R_1\theta_{x1} + R_2\theta_{x2}) \cos\psi - (R_1\theta_{y1} + R_2\theta_{y2}) \sin\psi
	\end{aligned}
	\label{eq:delta_aux}
\end{equation}

\begin{equation}
	\begin{aligned}
		\partial_{x_1}\delta &= (\sin\psi + \Delta_{\text{trans}}\partial_{x_1}\psi) \cos\beta_h + (\Delta_{\text{rot}}\partial_{x_1}\psi) \sin\beta_h \\
		\partial_{y_1}\delta &= (\cos\psi + \Delta_{\text{trans}}\partial_{y_1}\psi) \cos\beta_h + (\Delta_{\text{rot}}\partial_{y_1}\psi) \sin\beta_h \\
		\partial_{x_2}\delta &= (-\sin\psi + \Delta_{\text{trans}}\partial_{x_2}\psi) \cos\beta_h + (\Delta_{\text{rot}}\partial_{x_2}\psi) \sin\beta_h \\
		\partial_{y_2}\delta &= (-\cos\psi + \Delta_{\text{trans}}\partial_{y_2}\psi) \cos\beta_h + (\Delta_{\text{rot}}\partial_{y_2}\psi) \sin\beta_h
	\end{aligned}
	\label{eq:ddelta_xy}
\end{equation}

\begin{equation}
	\partial_{z_1}\delta = -\sin\beta_h, \quad \partial_{z_2}\delta = \sin\beta_h
	\label{eq:ddelta_z}
\end{equation}

\begin{equation}
	\partial_{\theta_{x1}}\delta = R_1 \sin\psi \sin\beta_h, \quad \partial_{\theta_{x2}}\delta = R_2 \sin\psi \sin\beta_h
	\label{eq:ddelta_thetax}
\end{equation}

\begin{equation}
	\partial_{\theta_{y1}}\delta = R_1 \cos\psi \sin\beta_h, \quad \partial_{\theta_{y2}}\delta = R_2 \cos\psi \sin\beta_h
	\label{eq:ddelta_thetay}
\end{equation}

\begin{equation}
	\partial_{\theta_{z1}}\delta = R_1 \cos\beta_h, \quad \partial_{\theta_{z2}}\delta = R_2 \cos\beta_h
	\label{eq:ddelta_thetaz}
\end{equation}

\text{Abordagem de Suavização Global (Smooth Operator):}

\begin{equation}
	g_1 = (\delta - b_t) \tanh(\sigma(\delta - b_t)), \quad g_2 = (\delta + b_t) \tanh(\sigma(\delta + b_t))
	\label{eq:smooth_g}
\end{equation}

\begin{equation}
	f(\delta, b_t) = \delta + \frac{1}{2}(g_1 - g_2)
	\label{eq:smooth_f}
\end{equation}

\begin{equation}
	\begin{aligned}
		\partial_{\delta} f &= 1 + \frac{1}{2} \bigg[ \tanh(\sigma(\delta - b_t)) + \sigma(\delta - b_t)(1 - \tanh^2(\sigma(\delta - b_t))) \\
		&\quad - \tanh(\sigma(\delta + b_t)) - \sigma(\delta + b_t)(1 - \tanh^2(\sigma(\delta + b_t))) \bigg] \\
		\partial_{b_t} f &= \frac{1}{2} \bigg[ -\tanh(\sigma(\delta - b_t)) - \sigma(\delta - b_t)(1 - \tanh^2(\sigma(\delta - b_t))) \\
		&\quad - \tanh(\sigma(\delta + b_t)) - \sigma(\delta + b_t)(1 - \tanh^2(\sigma(\delta + b_t))) \bigg]
	\end{aligned}
	\label{eq:smooth_df}
\end{equation}

\begin{equation}
	f_1(\delta, b_t) = \partial_{\delta} f \cdot \dot{\delta} + \partial_{b_t} f \cdot \dot{b}_t
	\label{eq:smooth_f1}
\end{equation}

\begin{equation}
	\partial_{q_i} f = \partial_{\delta} f \cdot \partial_{q_i}\delta + \partial_{b_t} f \cdot \partial_{q_i}b_t
	\label{eq:chain_rule_f}
\end{equation}